\documentclass[a4paper, 14pt]{extarticle}
\usepackage[T2A]{fontenc}
\usepackage[utf8]{inputenc}
\usepackage[russian]{babel}
\usepackage{geometry, amsmath, amssymb, mathtext}
\geometry{left=3cm, right=1.5cm, top=2cm, bottom=2cm}

\begin{document}
	
	\begin{titlepage}
		\centering
		{\small МИНИСТЕРСТВО НАУКИ И ВЫСШЕГО ОБРАЗОВАНИЯ \\ РОССИЙСКОЙ ФЕДЕРАЦИИ\\
		МОСКОВСКИЙ ФИЗИКО-ТЕХНИЧЕСКИЙ ИНСТИТУТ\\
		(НАЦИОНАЛЬНЫЙ ИССЛЕДОВАТЕЛЬСКИЙ УНИВЕРСИТЕТ)\\
		\vspace {1cm}
		Кафедра вакуумной электроники\\}
		
		\vspace{4cm}
		
		Лабораторная работа \\
		по курсу Вакуумная электроника\\
		{\Large \textbf{Напыление тонких плёнок в вакууме}\\}
		
		\vspace{3cm}
		
		\begin{flushright}
			Выполнил \\
			студент группы Б04-507 \\
			Рудачихин П.А. \\
			\vspace{0.5cm}
			Преподаватель: \\
			Пушкин А.С.
		\end{flushright}
		
		\vfill
		
		Долгопрудный \\
		2026
	\end{titlepage}
	
	\section{Аннотация}
	
	Цель данной работы -- ознакомление с методами термовакуумного напыления тонких плёнок. Производится термовакуумное напыление алюминиевой плёнки толщиной 10 нм на стеклянную пластинку, измеряется толщина плёнки оптическим методом.
	
	\section{Введение}
	
	\hspace{4mm} Рассмотрим простейшую модель термовакуумного напыления с резистивным испарителем (см. рис. 1). Испаряемое вещество находится в виде маленькой капли (точечный источник) непосредственно на испарителе, который нагревается пропускаемым через него током до температуры $T$. Над испарителем на расстоянии $R$ горизонтально расположена плоскость подложки.
	
	Для качественного формирования плёнки важно минимизировать появление частиц остаточных газов в пространстве между испарителем и подложкой, т.е. длина свободного пробега остаточных газов $\lambda$ должна быть много больше $R$. Предположим, что частицы остаточных газов с концентрацией $n$ и давлением $P$ имеют одинаковые радиус $r$ и температуру $T_0$, тогда
	
	\begin{equation}
		\lambda = \frac{1}{\sqrt{2}\cdot n\cdot 4\pi r^2} = \frac{kT_0}{4\sqrt{2}\pi P r^2}
	\end{equation}	
	
	Подставив средний радиус частиц воздуха $r = 0,187$ нм и комнатную температуру $T= 298$ К, получаем
	
	\begin{equation}
		\lambda \approx \frac{6,62\cdot 10^{-3}}{P}, \ \ \ \lambda(\text{см}) \approx \frac{8,71\cdot 10^{-3}}{P(\text{Торр})}
	\end{equation}	
	
	Оценим время образования монослоя остаточного газа на поверхности подложки. Выразим поток молекул на подложку через давление и температуру
	
	\begin{equation}
		j = \frac{n \langle v \rangle}{4} = \frac{P}{4kT_0} \sqrt{\frac{8kT_0}{\pi m}} = \frac{P}{\sqrt{2\pi mkT_0}} = \frac{PN_A}{\sqrt{2\pi MRT_0}}
	\end{equation}
	
	Подставим $M = 29$ г/моль и комнатную температуру
	
	\begin{equation}
		j = 2,12\cdot 10^{20} \cdot P(\text{Торр})
	\end{equation}
	
	Положив коэффициент прилипания равным единице, а поверхностную плотность атомов в монослое ???, найдём время образования монослоя
	
	\begin{equation}
		t_m = \frac{n_s}{j} = 
	\end{equation}
	
	
	Считая испарение изотропным, можно ввести массовый поток 
	
	\begin{equation}
		\mu := \frac{dm}{Sdt} = \frac{\dot{m}}{4\pi r^2}
	\end{equation} 
	
	где $r$ -- расстояние до источника.
	
	Из геометрических соображений (см. рис. 2) $r^2 = R^2 + x^2$, $dS_x = dS/\cos\theta$. Поток массы на плоскости
	
	\begin{equation}
		\mu_x := \frac{dm}{dS_x dt} = \mu \cos\theta = \frac{\dot{m}}{4\pi}\frac{R}{(R^2 + x^2)^{3/2}}
	\end{equation} 
	
	С другой стороны, приняв во внимание $dV = dh \cdot dS_x$, имеем
	
	\begin{equation}
		\mu_x := \frac{dm}{dS_x dt} = \frac{\rho dV}{dS_x dt} = \rho \dot{h}
	\end{equation} 
	
	Получаем уравнение
	
	\begin{equation}
		\dot{h} = \frac{\dot{m}}{4\pi\rho}\frac{x}{(R^2 + x^2)^{3/2}} \Rightarrow h = \frac{m}{4\pi\rho}\frac{R}{(R^2 + x^2)^{3/2}}
	\end{equation}
	
	Константы нет, так как на бесконечности h=0. В частности, в ближайшей к источнику точке плоскости
	
	\begin{equation}
		h_0 = \frac{m}{4\pi\rho R^2}
	\end{equation}
	
	
	 
	
	В данной работе используется резистивный тип испарителя, поэтому необходимо, чтобы температура плавления испарителя была выше температуры испаряемого вещества, давление паров материала испарителя было много меньше давления паров испаряемого вещества, материалы испарителя и испаряемого вещества не вступали в химическую реакцию. Испаряемое вещество должно смачивать испаритель, чтобы обеспечить хороший тепловой контакт и механическую стабильность.
	
	\section{Описание установки}
	
	
	
	\section{Результаты и обсуждения}
	
	
	
	\section{Выводы}
	
	
	
	\section{Приложение}
	
\end{document}
